\begin{abstract}
The widespread adoption of deep learning has profoundly influenced the analysis of diverse data types, including medical time-series modalities such as electroencephalography (EEG). Despite significant advancements, neural networks trained on EEG data frequently struggle with robustness and calibration when encountering realistic clinical conditions characterized by noise, electrode failures, and signal artifacts. Such vulnerabilities severely limit their reliability and clinical applicability, highlighting the need for structured and systematic robustness evaluation.

To bridge this critical gap, this work introduces TUSZ-C, a novel EEG dataset derived from the widely recognized Temple University Seizure Corpus (TUSZ). TUSZ-C systematically integrates clinically relevant corruptions, such as Gaussian noise, spectral masking, channel dropout, and signal dropout, to simulate realistic data perturbations encountered in clinical environments. This dataset enables comprehensive benchmarking of neural network robustness specifically tailored to EEG analysis.

This work further evaluates multiple neural network techniques known to enhance robustness and calibration, including Mixup, Manifold Mixup, Multi-input Multi-output (MIMO) ensembles, and Test-Time Adaptation (TTA).  All models are implemented using a ResNet-18 backbone and trained on the Temple University Seizure Corpus (TUSZ). On the clean test set, MIMO achieves the best trade-off, with an F1 score of 33.7\%, accuracy of 98.6\%, and Expected Calibration Error (ECE) of 0.07. Manifold Mixup follows closely, with an F1 score of 29.1\%, accuracy of 97.9\%, and ECE of 0.11. However, combining MIMO and Manifold Mixup substantially degrades performance, resulting in an F1 score of only 25.6\%, accuracy of 87.1\%, and elevated ECE of 0.15. All models experience some performance degradation on the TUSZ-C test set, which includes real-world signal corruptions; however, MIMO remains the most robust overall, achieving the best balance of accuracy and calibration under perturbations.

To further improve model reliability, we apply Test-Time Adaptation (TTA) by updating batch normalization statistics at inference. TTA reduces ECE consistently on both clean and corrupted test sets, particularly benefiting overconfident models. Our results demonstrate that combining lightweight training-time strategies with TTA offers a practical path toward robust, well-calibrated EEG classifiers for real-world deployment.
\end{abstract}